% Plantoir house style: tables in a handout (the skill's section 9). Three
% patterns to copy into a handout, each compiled and checked at 80 and 600 dpi:
%
%   A. A NEW table: nicematrix, every rule drawn after the heading's colour.
%   B. An EXISTING colortbl table that is to stay a tabular: \rowcolor with
%      its two overhang arguments (inset 0.5pt) - never \cellcolor.
%   C. Tables side by side: one top edge, one row pitch.
%
% Never mix the two kinds in one table. The preamble lines a handout needs
% are marked "Preamble" below; the starters carry them already, except
% [table] on xcolor, which only pattern B needs.
%
% WHY (measured 2026-10-31, tectonic, poppler's pdftoppm and CoreGraphics):
% - A shaded heading made with \cellcolor inside a \multicolumn and one made
%   with a bare \rowcolor draw byte-identical PDF: colortbl paints each
%   cell's fill AFTER the rule above it and the rule on its left. At 600 dpi
%   nothing is covered; at 80 dpi (the house check, and roughly what a
%   screen shows) poppler loses the whole top rule (58 px of it) and 19 px
%   of the left rule - with either command. A table with no fill loses none.
%   So moving from \cellcolor to a bare \rowcolor fixes nothing.
% - What fixed the teacher's tables was \rowcolor WITH the two overhang
%   arguments below (the fill stops 0.5pt short of each side rule: 0 px of
%   the side rules lost) plus a global redraw of every tabular's outer frame
%   (0 px lost anywhere). Without the redraw, B's top rule depends on where
%   the table sits: set at 12 sub-pixel page positions (2026-11-01), 54-56
%   of its 58-59 px were lost at 80 dpi at 6 of the 12, exactly where a
%   bare \rowcolor lost its top rule too, and 0 at 600 dpi at all 12.
% - nicematrix (pattern A) paints the colours FIRST and every rule after:
%   0 px lost at 80, 100, 150 and 600 dpi, with no redraw, and the same
%   geometry as a tabular. That is why it is the pattern for new tables.
% - Side by side, plain tabulars with a two-line heading in one of them
%   drifted 46 px at 600 dpi; pattern C put every rule of three tables at
%   the same height.
\documentclass[11pt]{article}
\usepackage[letterpaper, margin=0.75in, bottom=0.9in]{geometry}
\usepackage{lmodern}
\usepackage[T1]{fontenc}
\usepackage{amsmath}
\usepackage[table]{xcolor}  % Preamble, pattern B only: [table] loads colortbl
\usepackage{array}          % Preamble: r<{...} in a column spec
\usepackage{nicematrix}     % Preamble: patterns A and C
\definecolor{tablehead}{RGB}{225,232,245}             % Preamble: the heading row's shade
\newcommand{\headstrut}{\rule[-1.1em]{0pt}{2.6em}}    % Preamble: one heading-row height
\setlength{\parindent}{0pt}
\pagestyle{empty}
\providecommand{\patternbreak}{\par\bigskip\bigskip}

\begin{document}

% ---------- A. A new table ----------
% Per column, by its values: x (1 to 4) has no negative value, so it is
% centred (c); y holds -3, -4, -3 and 0, so it is right-aligned with 0.25cm
% of padding before the rule. Zero is not negative and does not make a column
% "all the same sign". The heading over a right-aligned column is centred
% with \multicolumn{1}{c}{...}; in a NiceTabular the [hvlines] option draws
% every rule, so the heading's spec needs no | and no \hline is written.
{\renewcommand{\arraystretch}{1.25}%
\begin{NiceTabular}{cr<{\hspace{0.25cm}}}[hvlines]
\CodeBefore
  \rowcolor{tablehead}{1}
\Body
$x$ & \multicolumn{1}{c}{$y$} \\
$1$ & $-3$ \\
$2$ & $-4$ \\
$3$ & $-3$ \\
$4$ & $0$ \\
\end{NiceTabular}}

\patternbreak

% ---------- B. An existing colortbl table ----------
% Keep the full border: | at both ends of the spec, \hline above the heading
% row and below the last row. Shade the heading with \rowcolor at the START
% of the row, with both overhang arguments, and no \cellcolor in any cell.
% Its top rule VANISHES at 80 dpi at about half of page positions, and 600
% dpi never shows it: look at the top edge in the 80 dpi pages, and if it is
% gone rewrite the table as pattern A - unless the preamble already redraws
% every tabular's outer frame last (a teacher's template may carry one, as
% an "Outer border, drawn last" block), which closes it. A table being rewritten anyway
% becomes pattern A. In a plain tabular the heading's \multicolumn repeats the rules:
% {|c|} over the first column, {c|} over the others.
{\renewcommand{\arraystretch}{1.25}%
\begin{tabular}{|c|r<{\hspace{0.25cm}}|}
\hline
\rowcolor{tablehead}[\dimexpr\tabcolsep-0.5pt\relax][\dimexpr\tabcolsep-0.5pt\relax]
\multicolumn{1}{|c|}{$x$} & \multicolumn{1}{c|}{$y$} \\
\hline
$1$ & $-3$ \\ \hline
$2$ & $-4$ \\ \hline
$3$ & $-3$ \\ \hline
$4$ & $0$ \\ \hline
\end{tabular}}

\patternbreak

% ---------- C. Tables side by side ----------
% Every table in one row: a [t] minipage starting with \vspace{0pt}, [t] on
% the table, ONE \arraystretch for the row, and \headstrut in every heading
% row so a one-line heading is as tall as a two-line one. Where tables list
% the same x values, the rows then line up across the page. In a heading
% cell made with \multicolumn the strut goes INSIDE it ("Misplaced \omit"
% otherwise). A two-line heading is \Block{}{First\\differences}.
% Decimals: in a column mixing whole numbers and decimals, every value has
% the same places (2.0 beside 4.5, 0.0 beside 0.5) - never a bare whole
% number, never \phantom. A column of whole numbers gets no decimals.
{\renewcommand{\arraystretch}{1.25}%
\begin{minipage}[t]{0.3\linewidth}\vspace{0pt}\centering
\begin{NiceTabular}[t]{r<{\hspace{0.25cm}}r<{\hspace{0.25cm}}}[hvlines]
\CodeBefore \rowcolor{tablehead}{1} \Body
\multicolumn{1}{c}{\headstrut$x$} & \multicolumn{1}{c}{$y$} \\
$-2$ & $-6$ \\
$-1$ & $-10$ \\
$0$ & $-14$ \\
$1$ & $-18$ \\
\end{NiceTabular}
\end{minipage}\hfill
\begin{minipage}[t]{0.3\linewidth}\vspace{0pt}\centering
\begin{NiceTabular}[t]{ccc}[hvlines]
\CodeBefore \rowcolor{tablehead}{1} \Body
\headstrut$x$ & $y$ & \Block{}{First\\differences} \\
$0$ & $1$ & \\
$1$ & $3$ & \\
$2$ & $5$ & \\
$3$ & $7$ & \\
\end{NiceTabular}
\end{minipage}\hfill
\begin{minipage}[t]{0.3\linewidth}\vspace{0pt}\centering
\begin{NiceTabular}[t]{cc}[hvlines]
\CodeBefore \rowcolor{tablehead}{1} \Body
\headstrut$t$ (s) & $h$ (m) \\
$0$ & $2.0$ \\
$1$ & $4.5$ \\
$2$ & $0.5$ \\
$3$ & $0.0$ \\
\end{NiceTabular}
\end{minipage}}

\end{document}
